% Generated draft: requires mathematical and source review.
\documentclass[11pt]{article}
\usepackage[margin=25mm]{geometry}
\usepackage{fontspec}
\setmainfont{Latin Modern Roman}
\usepackage{amsmath,amssymb,mathtools,unicode-math}
\setmathfont{Latin Modern Math}
\usepackage{enumitem,needspace,xurl,hyperref}
\setlength{\parindent}{0pt}
\setlength{\parskip}{5pt}
\setlength{\emergencystretch}{3em}
\newcommand{\N}{\mathbb{N}}
\newcommand{\Z}{\mathbb{Z}}
\newcommand{\Q}{\mathbb{Q}}
\newcommand{\R}{\mathbb{R}}
\newcommand{\C}{\mathbb{C}}
\newcommand{\F}{\mathbb{F}}
\newcommand{\A}{\mathbb{A}}
\newcommand{\PP}{\mathbb{P}}
\newcommand{\E}{\mathbb{E}}
\newcommand{\eps}{\varepsilon}
\newcommand{\dd}{\,\mathrm{d}}
\begin{document}
\section*{UnsolvedMath ID: 30005416}
\textbf{Source problem code:} catalogue entry
\subsection{Definitions and mathematical statement}
Let \(q\) be a formal variable, and define
\[
P(q)=\frac{1}{(1-q)^3}\prod_{i=2}^{\infty}\frac{1}{(1-q^i)^2}\in\Z[[q]].
\]
The product is interpreted formally: each coefficient depends on only finitely many factors. Write \(\operatorname{AHP}_{X,O}(q)\) for the arc Hilbert--Poincar\'e series of a pointed singularity \((X,O)\), as named in the supplied source. The record does not define this invariant or specify the algebra, grading, or filtration used to construct it; consequently, its precise definition remains to be supplied from the underlying report.

The source identifies ADE surface singularities with rational double point surface singularities. It reports, as Theorem 1 attributed to its reference [7], that every such singularity satisfies
\[
\operatorname{AHP}_{X,O}(q)=P(q).
\]
\textbf{Question, with unresolved scope.} Does equality with this series characterize the class of ADE surface singularities? More explicitly, once the intended class \(\mathcal C\) of pointed surface singularities and the definition of \(\operatorname{AHP}\) have been fixed, is the following implication true?
\[
\forall (X,O)\in\mathcal C,\qquad
\operatorname{AHP}_{X,O}(q)=P(q)
\ \Longrightarrow\
(X,O)\text{ is an ADE surface singularity}.
\]
Together with the reported forward implication, this would give an equivalence on \(\mathcal C\). The record does not specify \(\mathcal C\), the ground field or its characteristic, the algebraic, analytic, or formal category, or hypotheses such as reducedness, normality, or isolation of the singularity. This is therefore a formulation schema requiring those qualifications, rather than a fully specified universally quantified conjecture. Because the reported series is the same for every ADE singularity, the proposed characterization concerns membership in the ADE class; the displayed formula cannot distinguish individual ADE types.

For context, the source says that this invariant detects regularity and contains ingredients relevant to singularity theory, referring to its reference [8], and that its computation is generally difficult. It also reports, as Theorem 3 attributed to its reference [5], that for the dual-number scheme
\[
X=\operatorname{Spec}\bigl(K[x]/(x^2)\bigr),\qquad O=(x),
\]
one has
\[
\operatorname{AHP}_{X,O}(q)
=\prod_{\substack{i\geq 1\\ i\equiv 1,4\pmod 5}}\frac{1}{1-q^i}.
\]
Here \(K\) is the unspecified field appearing in the source; no characteristic hypothesis is supplied. This contextual example is not a surface singularity and does not specify the comparison class in the question. All reported formulas and regularity claims require checking against the original definitions and hypotheses.
\subsection{Short English statement}
The question asks whether a surface singularity whose arc Hilbert--Poincar\'e series equals the displayed common ADE series must be a rational double point. The source reports the forward implication. The invariant's definition, the ground field, and the class of singularities under comparison are missing, so the converse remains a statement schema requiring human review. The supplied open status is dated curation, not an independently verified assessment.
\subsection{Sources}
\begin{itemize}
\item Supplied local UnsolvedMath registry record, permanent ID \(30005416\), source problem code \texttt{catalogue entry}, titled ``Arc Hilbert--Poincar\'e Series of ADE Surface Singularities.'' No individual proposer is identified; the supplied proposed year is 2023.
\item \emph{Resolutions in Local Algebra and Singularity Theory} (2023), Oberwolfach Reports, pp.\ 327--395, DOI \texttt{10.4171/owr/2023/6}. Supplied report citation and literature-source link labeled ``original OWR question'': \url{https://doi.org/10.4171/owr/2023/6}.
\item Supplied literature source labeled ``Mourtada, jet schemes of rational double points'': \url{https://webusers.imj-prg.fr/~hussein.mourtada/PointDouble30Sept.pdf}. The record does not establish how this source defines or supports the particular arc Hilbert--Poincar\'e series used here.
\item The excerpt attributes the rational-double-point formula to Theorem 1, reference [7]; the dual-number formula to Theorem 3, reference [5]; and background on the invariant to reference [8]. Bibliographic entries for [5], [7], and [8] are not supplied, and their identities cannot be recovered reliably from this record alone.
\item Literature assessment dated exactly \texttt{2026-08-22}: the supplied curation labels the problem open, states that the series has been computed for rational double points, and describes the characterization question as unresolved. This is supplied dated curation, not independently verified here. The background explicitly requests independent verification of time-sensitive status and cited claims.
\item The registry labels the statement \texttt{corrected\_verified} and records that the official report was checked directly. These are supplied provenance claims; no independent source visit or mathematical verification is asserted in this draft.
\end{itemize}
\end{document}
